\documentclass[11pt]{article}
\usepackage[a4paper,margin=27mm]{geometry}
\usepackage{amsmath,amssymb,amsthm,hyperref}
\newcommand{\PSH}{\operatorname{PSH}}
\newcommand{\FS}{\operatorname{FS}}
\newcommand{\vol}{\operatorname{vol}}
\newcommand{\ord}{\operatorname{ord}}
\newtheorem{theorem}{Theorem}
\newtheorem{lemma}[theorem]{Lemma}
\title{The entire genuine big surface cone through its Zariski negative part}
\author{}
\date{6 October 2026}
\begin{document}
\maketitle
\begin{abstract}
For a smooth projective surface, this candidate identifies the entire
genuine numerical plurisubharmonic cone of a big rational bundle with
the nef cone of its Zariski positive part translated by the canonical
negative-part potential. The proof retains actual sections, arbitrary
directed nets, and all analytic points. It derives nef closure from the
exact actual-ample theorem and constructs the negative-part limit
without general nonample compactness. It would repair the full big
surface closure gate, but has not undergone independent review and
does not solve original P16.5 in arbitrary dimension or construct its
entire finite-energy equation.
\end{abstract}

\section{Exact genuine objects and the sole analytic premise}
Let $k$ be any algebraically closed trivially valued field of
characteristic zero, $X$ a smooth integral projective surface over $k$,
and $H$ an actual ample rational bundle. All functions are on the
original compact Hausdorff space $X^{\rm an}$; $x_0$ denotes the
generic trivial valuation. A rational bundle is an actual element of
$\operatorname{Pic}(X)\otimes\mathbf Q$.
For an actual rational $B$, a generically finite FS potential is
\[
 f=\frac1m\max_j(\log|s_j|+\lambda_j),\qquad
 0\ne s_j\in H^0(X,mB),\quad\lambda_j\in\mathbf R.
                                                               \tag{1}
\]
Its canonical trivial section norm is used; poles are allowed.
The genuine numerical cone consists exactly of the functions
\[
 \PSH(\theta)=\{u=\inf_i f_i:\ f_i\searrow\text{ at all points},
       \ f_i\in\FS^{\rm gf}(B_i),\ c_1(B_i)\to\theta,
       \ u(x_0)>-\infty\}.                                \tag{2}
\]
The index set is an arbitrary directed set. This is a one-net
definition, not an iterated or divisorial closure. Its normalized
part has $u(x_0)=0$.

The analytic premise $G_0$ is the exact independently supplied
actual-ample theorem on this original $X,k$: the genuine normalized
ample cone is compact in divisorial coordinates, is determined at
all points by those coordinates, admits decreasing regularization
by finite continuous actual FS functions on the same bundle,
and is decreasing closed. Its whole mathematical body is retained
in the research record; no general-big theorem inside a historical
discussion is granted. The earlier research132 nef statements are
not premises; their necessary structural proof is reconstructed below.

We retain the ordinary algebraic facts that the ample numerical cone
is open, nef plus ample is ample, the nef cone is closed, and
divisible ample powers are generated. Smooth projective surface
intersection numbers, Serre vanishing, surface Riemann--Roch and
Serre duality are the other background facts. No nonample psh
closure theorem is included in them.

\section{Section embeddings and the entire nef structural theorem}
An actual ample rational bundle $C$ has a generated power with
weight-zero FS potential identically zero: at each valuation-ring
center some generator is a unit of norm one, and every generator
has norm at most one. Tensoring with this system makes (1) the same
all-point function on $B+C$. Tensor products of all section monomials,
after clearing denominators, realize sums and positive integer
powers of FS potentials exactly.

Consequently a genuine function in class $\theta$ embeds unchanged
in every actual ample bundle $A$ with $c_1(A)-\theta$ ample. Indeed
on a tail $A-B_i$ is actual ample, and the zero system tensors each
defining stage into $A$. This is already one decreasing net.
It proves finite divisorial evaluations and all-point determination
for any two already genuine functions in one numerical class,
by putting them into one sufficiently large actual ample cone.
It does not prove membership for a new nonample limit.

\begin{lemma}[The entire nef cone]
For every actual nef rational bundle $P$,
\[
 \PSH(c_1(P))=\bigcap_{\epsilon\in\mathbf Q_{>0}}
                      \PSH(P+\epsilon H).                  \tag{3}
\]
It has one-net decreasing regularization by finite continuous FS
potentials with numerical classes tending to $c_1(P)$, is closed
under all decreasing nets with finite limiting generic value, and
its normalized part is compact in the divisorial topology.
\end{lemma}
\begin{proof}
Forward membership follows by the section embedding, since
$P+\epsilon H$ is actual ample. For reverse membership, let $u$
belong to every ample cone on the right and choose its finite
continuous FS regularizers $f_{\epsilon,j}\searrow u$ there.
Consider triples $(\epsilon,j,a)$ with $a>0$ and order them by
decrease of $\epsilon$, of $a$, and of the function
$f_{\epsilon,j}+a$. They form a directed preorder. To see this,
choose $\epsilon_3$ below two given values and $a_3$ smaller than
their positive shifts. The finite continuous function
$R=\min(f_{\epsilon_1,j_1}+a_1,f_{\epsilon_2,j_2}+a_2)$
strictly exceeds $u+a_3$ even where $u=-\infty$.
The open negative sets of
$\exp(f_{\epsilon_3,j}+a_3)-\exp R$ cover the compact carrier.
A finite subcover and directedness give one $j_3$ with strict
comparison everywhere. The shifted system is FS on the actual
$P+\epsilon_3H$. Its numerical classes tend to $c_1(P)$, and
finite thresholds at each point show that its decreasing limit
is exactly $u$. This proves (3) with one defining net.

If $u_i\searrow u$ in the nef cone, fixed ample $P+\epsilon H$
contains the same functions; $G_0$ puts their all-point limit
there when the limiting generic value is finite. Apply (3).
For normalized compactness, each compact normalized ample cone
for $P+\epsilon H$, $0<\epsilon\leq1$, embeds in that for
$P+H$. Its image is closed in this Hausdorff coordinate space.
Their intersection is compact. Representatives with the same
divisorial coordinates coincide all-pointwise inside the one ample
cone by $G_0$. Equation (3) identifies this intersection as the
entire normalized nef cone, with its original representatives.
\end{proof}
The construction also permits passage to varying nef classes:
if $P_j$ are actual nef rational bundles with $c_1(P_j)\to c_1(P)$,
normalized $v_j\in\PSH(c_1(P_j))$, and a subnet converges in one
dominating normalized ample cone, its limit belongs to
$\PSH(c_1(P))$. For fixed rational $\delta>0$, eventually
$P+\delta H-P_j$ is actual ample; the section embedding puts
$v_j$ in the compact closed ample cone of $P+\delta H$.
The limit belongs to every such cone, and (3) applies. This argument
uses compact subnets, not sequential compactness or a countability
assumption on $k$.

\section{Algebraic negative parts with their complete finite proof}
Let $L$ be an actual big rational bundle. Choose an effective rational
divisor $D$ representing it, by taking a nonzero section of a
divisible power. Here is a complete finite construction of its
Zariski decomposition, in the spirit of Bauer's maximal nef argument.

Write $D=\sum d_iC_i$ with finitely many distinct prime curves.
The real polytope
\[
 K=\{0\leq p_i\leq d_i:\ \sum_i p_i(C_i\cdot C_j)\geq0
                                      \text{ for every }j\}
\]
is nonempty and rational. Its vectors describe exactly the effective
nef subdivisors: a curve outside the displayed support has
nonnegative intersection with all its components. Maximize
$\sum p_i$ and choose a rational maximizing vertex. Coordinatewise
maxima preserve these inequalities, since off-diagonal curve
intersections are nonnegative. Thus the maximizer is the unique
greatest feasible vector. Let $P=\sum p_iC_i$ and $N=D-P$.
If $C_j$ occurs in $N$, then $P\cdot C_j=0$: otherwise adding
a sufficiently small positive multiple of $C_j$ preserves every
inequality and the cuboid, contradicting maximality.

The intersection matrix $Q$ on the support of $N$ is negative
definite. If it were not, its largest eigenvalue would be nonnegative
and admit a nonnegative eigenvector $q\ne0$. Indeed a unit vector
maximizing the Rayleigh quotient can be replaced by its coordinatewise
absolute value without decreasing that quotient, because all
off-diagonal entries are nonnegative. The maximizer is an eigenvector.
Then $\sum q_iC_i$ has nonnegative intersection with every supported
curve and every other curve, so is a nonzero effective real nef
divisor. Adding a sufficiently small multiple to $P$ stays inside
the real feasible polytope and improves the maximum, a contradiction.
This proves existence of rational $P,N$, with $P$ nef, $N\geq0$,
negative definite support, and $P\cdot C_i=0$ on that support.

We need its exact minimality and section implications, which are
proved rather than silently imported. For a negative definite
matrix $Q$ with nonnegative off-diagonal entries, $Q^{-1}$ has
nonpositive entries. To prove this put $A=-Q$ and solve $Ax=b\geq0$.
If $x=x_+-x_-$ with disjoint nonnegative parts and $x_-\ne0$, then
$0\leq b\cdot x_-=x_+^TAx_--x_-^TAx_-<0$:
the first term is nonpositive and the second strictly positive.
Hence $x\geq0$. Applying this to the coordinate vectors proves
$A^{-1}\geq0$, as required.

If numerically $L=Q_0+M$ with $Q_0$ nef and $M\geq0$, split
$M=\sum m_iC_i+M'$ over the negative support, where $M'$ has no
such components. With $a_i$ the coefficients of $N$, nefness gives
$Q(a-m)_i-M'\cdot C_i\geq0$.
The inverse-sign result forces $a-m\leq0$, so $M\geq N$.
Thus the negative part is coefficientwise minimal in every
nef-plus-effective decomposition. It is unique and depends only
on the numerical class of $L$; comparison of two decompositions
proves this directly. The actual positive bundle is $P=L-N$.

If $T\in|mL|$ is effective and $mN$ integral, write
$T=\sum t_iC_i+T'$. Orthogonality yields
$Q(t-ma)_i=-T'\cdot C_i\leq0$.
Again $Q^{-1}\leq0$ gives $t\geq ma$. Every actual section of
$mL$ is therefore divisible by the canonical section of $mN$.
Conversely multiplication by that section embeds $H^0(mP)$ in
$H^0(mL)$. This proves the actual isomorphism
\[
 H^0(X,mP)\xrightarrow{\ \cdot s_{mN}\ }H^0(X,mL)
                   \quad\text{for all sufficiently divisible }m.
                                                               \tag{4}
\]

For $L_\epsilon=L+\epsilon H$ let $N_\epsilon$ be its negative
part and $P_\epsilon=L_\epsilon-N_\epsilon$.
Minimality, applied to $P+\epsilon H+N$, gives
$0\leq N_\epsilon\leq N$. Applying it again when
$0<\epsilon'<\epsilon$ gives $N_\epsilon\leq N_{\epsilon'}$.
These divisors are supported on the one finite support of $N$,
so they converge coefficientwise to $N_*\leq N$ as
$\epsilon\downarrow0$. Closedness of the nef cone gives
$L-N_*$ nef. Minimality for $L$ forces $N_*\geq N$.
Consequently
\[
 N_\epsilon\uparrow N,\qquad c_1(P_\epsilon)\to c_1(P).
                                                               \tag{5}
\]
This proves the required continuity directly, without assuming
local finiteness of chambers or finite generation of the moving ring.

\section{The entire big surface cone, with one-net membership}
For an effective rational divisor $N$ let
$\sigma_N=\log|s_{mN}|/m$, using any denominator-clearing $m$.
It is an actual one-section FS potential of $N$, independent of
the choice of $m$, normalized at $x_0$.

\begin{theorem}[Candidate structural repair]
For every actual big rational $L$ on the retained surface,
\[
 \PSH(c_1(L))=\{\sigma_N+v:v\in\PSH(c_1(P))\}.            \tag{6}
\]
This equality concerns all actual all-point functions in the entire
genuine cone. The translation is a homeomorphism on normalized
divisorial cones. In particular the entire big normalized cone is
compact and the genuine cone is closed under arbitrary decreasing
nets with finite limiting generic value.
\end{theorem}
\begin{proof}
Forward construction from the right is immediate but retains the
actual ideals: the nef lemma gives finite FS $f_i\searrow v$ on
$P+\epsilon_iH$. Tensor with $s_{mN}$, clearing denominators by
all section products. The resulting functions
$\sigma_N+f_i$ are actual FS on $L+\epsilon_iH$, decrease
at every point to $\sigma_N+v$, and have finite generic limit.
Their base ideal includes the exact normalized power of the
canonical negative-part ideal. Thus (2) proves genuine membership.

Conversely let $u\in\PSH(c_1(L))$, normalized at $x_0$.
Fix a rational $\epsilon>0$. From its original defining net (2),
the actual differences $L_\epsilon-B_i$ are ample on a tail.
Tensor their zero FS systems to place the unchanged defining
functions on the same actual $L_\epsilon$. Clear at each stage
the denominator of $N_\epsilon$ by all products of its sections.
Equation (4), now for $L_\epsilon$, divides them exactly by
$s_{mN_\epsilon}$. This gives actual FS functions $g_i$ on
$P_\epsilon$, with
$f_i=\sigma_{N_\epsilon}+g_i$ on every point where the first
factor is finite.

They decrease at every divisorial valuation, where
$\sigma_{N_\epsilon}$ is finite. Put the $g_i$ in one actual
ample cone by the section embedding. The maximum of these two
actual FS functions is actual FS on that same ample bundle:
clear their denominators, take all powers of each section system,
and take their union with its weights. Divisorial order and all-point
determination promote this decrease to all analytic points.
The nef closure lemma supplies $v_\epsilon\in\PSH(c_1(P_\epsilon))$
with one actual defining net. Its divisor coordinates are
\[
 v_\epsilon(z)=u(z)-\sigma_{N_\epsilon}(z)
               \quad\text{for every divisorial }z.         \tag{7}
\]
Its generic value is zero. The already constructed sum
$\sigma_{N_\epsilon}+v_\epsilon$ is genuine in $c_1(L_\epsilon)$.
The function $u$ is also genuine there by its unchanged original
net. Put both in one actual ample cone; (7) and determination give
their equality on all points. No subtraction at shared poles is
being used to define $v_\epsilon$.

Take rational $\epsilon\downarrow0$. The normalized $v_\epsilon$
all embed on a tail in one sufficiently large normalized ample
cone, since $P_\epsilon\to P$. This cone is compact, so there
is a convergent subnet. At every divisor (5) and (7) prescribe its
finite limit $u(z)-\sigma_N(z)$. The varying-nef-class argument
after (3) puts its actual all-point representative $v$ in
$\PSH(c_1(P))$. Its normalized generic value is zero.
The sum $\sigma_N+v$ is genuine in $c_1(L)$ by the forward
construction and equals $u$ at all divisor coordinates.
Embedding both in one actual ample cone proves equality on all
analytic points, including all poles. This proves (6).

On divisorial coordinates this map is translation by the fixed
finite values of $\sigma_N$, and is injective by ample
determination for the nef source. It is therefore a homeomorphism
in the stated coordinate topologies; normalized compactness follows.
If $u_i=\sigma_N+v_i$ decreases with finite limiting generic
value, the $v_i$ decrease divisorially. In a common ample cone
the maximum of two of them is genuine: take its two finite FS
regularizing nets on that same bundle, form their product-directed
pointwise maxima as above, and pass to their decreasing limit.
Divisorial determination promotes their order to all points.
Nef closure produces
$v=\inf_i v_i$; literal addition by $\sigma_N$ gives
$\inf_i u_i=\sigma_N+v$. At poles of $\sigma_N$ both sides are
$-\infty$; elsewhere its value is finite and commutes with the
infimum. The forward construction supplies one actual defining net.
\end{proof}

\section{Raw volume and the precise remaining scope}
The raw algebraic volume is $\vol(L)=\vol(P)=P^2$.
The first equality follows from (4) on divisible powers.
For completeness, nefness gives the upper bound
$\vol(P)\leq(P+\delta H)^2$ by multiplying by a nonzero section
of a divisible $\delta H$ and using ample Serre vanishing and
surface Riemann--Roch. Let $\delta\downarrow0$.
Since $L$ is big, (4) and this bound imply $P^2>0$.
For large divisible $m$, $H^0(K_X-mP)=0$, because its
intersection with $H$ is negative. Serre duality and
Riemann--Roch then give
$h^0(mP)\geq\chi(mP)=m^2P^2/2+O(m)$, proving the reverse
bound. This fixes the intended raw mass; it does not construct
an unbounded original-carrier operator by choosing an inverse.

The theorem repairs only the structural closure and compactness
gate on the entire big surface domain, relative to the exact
actual-ample analytic premise and the explicitly proved algebraic
steps. The full original P16.5 permits arbitrary dimension and
requires the entire relative finite-energy operator, continuous
directional derivatives, intrinsic targets, weighted limits,
comparison, variational attainment and all-solution uniqueness.
None of those remaining full-domain conclusions is asserted here.
In higher dimension a finite canonical divisor whose removal makes
the whole cone nef need not exist; a generalization needs its own
genuine geometric construction, rather than a formal replacement
of $N$ by an unspecified birational object.

This is an unreviewed root candidate. The negative-part proof was
compared with the complete four-page primary text of T. Bauer,
\emph{A simple proof for the existence of Zariski decompositions
on surfaces},
\href{https://arxiv.org/pdf/0712.1576v1}{arXiv:0712.1576v1}.
The polytope construction and its off-diagonal sign facts are
written in full; the matrix inverse, perturbation limit,
actual section division and entire genuine-cone transfer are
also proved here. No abstract extended-real flattening lemma,
general nonample envelope theorem or historical verdict is a premise.
The full actual-ample $G_0$ body and the unreviewed completed132
research paper are preserved separately for independent audit.
No additional original atom or scoped paper is counted as solved.
\end{document}
